\section{A1 for \texorpdfstring{$\Spin(2n)$}{Spin(2n)}: explicit one-plaquette coefficient bounds}
\label{app:A1_Spin_even}

This appendix proves \Cref{prop:A1_Spin_even}.
As in Appendix~\ref{app:A1_Spin_odd}, the Wilson weight factors through $\SO(2n)$,
so only representations with $\pi(-\mathbf 1)=\mathbf 1$ contribute.
We therefore work on $\SO(2n)$ and restrict to dominant integral highest weights (the ``tensor'' sector).

\subsection{Maximal torus and Wilson weight}
\label{subsec:A1_Spin_even_setup}

Fix $n\ge 4$ and set $G=\Spin(2n)$.
On $\SO(2n)$, every conjugacy class meets the maximal torus with eigenvalues
\begin{equation}
\label{eq:SO_even_eigs}
\{e^{\pm i\theta_1},\dots,e^{\pm i\theta_n}\},\qquad \theta_j\in[0,\pi].
\end{equation}
For the defining $2n$-dimensional representation $\rho$,
\begin{equation}
\label{eq:SO_even_trace}
\Re\Tr_{\rho}(U)=2\sum_{j=1}^n \cos\theta_j.
\end{equation}
Hence, up to normalization,
\Cref{eq:one_plaq_measure_general} corresponds on the torus to the symbol
\begin{equation}
\label{eq:SO_even_symbol}
W_t(\theta):=\exp\Big(t\sum_{j=1}^n \cos\theta_j\Big),
\qquad t:=\frac{\alpha\beta}{n}.
\end{equation}

\subsection{Determinantal coefficient formula}
\label{subsec:A1_Spin_even_determinant}

Irreducible representations of $\SO(2n)$ in the tensor sector are indexed by dominant integral highest weights.
For the present A1 gate it suffices to label them by partitions $\lambda=(\lambda_1\ge\cdots\ge\lambda_n\ge 0)$;
representations on which $-\mathbf 1\in\Spin(2n)$ acts nontrivially (spinor sector) do not appear in the Wilson expansion.

Define the type-$D_n$ shifts
\begin{equation}
\label{eq:SO_even_bi}
 b_i:=n-i\qquad(1\le i\le n),
\end{equation}
so that the Weyl vector has components $\rho_i=b_i$.
Define
\begin{equation}
\label{eq:SO_even_Ai}
A_i(\lambda):=\lambda_i+b_i=\lambda_i+n-i\in\N_0.
\end{equation}

We record the determinant-form Weyl character and Weyl integration formulas for the tensor sector of $\SO(2n)$,
since these are the only group-specific inputs needed to reduce the one-plaquette coefficients to moment determinants
(and they avoid any reliance on Toeplitz$+$Hankel machinery).

\begin{lemma}[Weyl character formula for $\SO(2n)$ in cosine-determinant form (tensor sector)]
\label{lem:SOeven_character_cosine}
Let $U\in T$ have eigenangles $\theta=(\theta_1,\dots,\theta_n)$ as in \eqref{eq:SO_even_eigs}.
For a tensor-sector irreducible representation indexed by a partition $\lambda$ of length at most $n$,
\begin{equation}
\label{eq:SOeven_character_formula}
\chi_{\lambda}(\theta)
=\frac{\det\big(\cos(A_i(\lambda)\,\theta_j)\big)_{1\le i,j\le n}}
        {\det\big(\cos(b_i\,\theta_j)\big)_{1\le i,j\le n}},
\qquad b_i=n-i,\quad A_i(\lambda)=\lambda_i+b_i.
\end{equation}
\end{lemma}

\begin{lemma}[Weyl integration formula in determinant form for $\SO(2n)$]
\label{lem:SOeven_weyl_integration}
There exists a constant $c_{D}(n)>0$ such that for every continuous class function $F$ on $\SO(2n)$,
\begin{equation}
\label{eq:SOeven_weyl_integration}
\int_{\SO(2n)}F(U)\,d\mu_{\mathrm{Haar}}(U)
= c_{D}(n)\int_{[0,\pi]^n}F(\theta)\,
\det\big(\cos(b_i\,\theta_j)\big)_{1\le i,j\le n}^{\,2}\,d\theta_1\cdots d\theta_n,
\qquad b_i=n-i.
\end{equation}
\end{lemma}

\begin{proof}
This is the Weyl integration formula for type $D_n$ written using the determinant presentation of the Weyl denominator.
Concretely, the type-$D_n$ Weyl denominator is
\(
\prod_{1\le i<j\le n}2\sin\big(\frac{\theta_i-\theta_j}{2}\big)\,2\sin\big(\frac{\theta_i+\theta_j}{2}\big)
\),
and the identity
\(
2\sin\big(\frac{\theta_i-\theta_j}{2}\big)\,2\sin\big(\frac{\theta_i+\theta_j}{2}\big)
= \cos\theta_j-\cos\theta_i
\)
shows that (up to a constant) it equals the Vandermonde

\paragraph{Denominator-minor convention (TN ratios).}
At several points below we form Vandermonde-normalized minor ratios (or equivalent determinant ratios) and invoke the TN one-row shift tool from Appendix~\ref{app:TP_tools}.
Whenever a would-be denominator minor can vanish, we apply the zero-case clause from Lemma~J.5:
if the pivot entry (hence the denominator minor) vanishes, then the corresponding numerator determinant also vanishes and the desired inequality is trivial; otherwise we work in the strictly-positive regime and proceed with the ratio bound.


 determinant in the variables $\cos\theta_j$.
Since $\cos(b\,\theta)$ is a degree-$b$ polynomial in $\cos\theta$ with nonzero leading coefficient, the determinant
$\det(\cos(b_i\theta_j))$ is a nonzero constant multiple of this Vandermonde, and squaring yields \eqref{eq:SOeven_weyl_integration}.
\end{proof}

\begin{proposition}[Determinantal formula for one-plaquette coefficients on $\SO(2n)$]
\label{prop:SOeven_det_formula}
Let $\lambda$ be a partition of length at most $n$.
Define the cosine moment kernel
\begin{equation}
\label{eq:SOeven_kernel_def}
M_t^{D}(p,q)
:=\frac{2}{\pi}\int_0^{\pi} e^{t\cos\theta}\,\cos(p\theta)\,\cos(q\theta)\,d\theta,
\qquad p,q\in\N_0.
\end{equation}
Let
\begin{equation}
\label{eq:SOeven_Dlambda_def}
D_{\lambda}^{D}(t)
:=\det\Big(M_t^{D}\big(A_i(\lambda),\,b_j\big)\Big)_{1\le i,j\le n}.
\end{equation}
Then with $t=\alpha\beta/n$,
\begin{equation}
\label{eq:SOeven_coeff_det_ratio}
 a_{\lambda}(\beta,\alpha)
=\E_{\nu_{\beta,\alpha}}\Big[\frac{\chi_{\lambda}(U)}{\dim\lambda}\Big]
=\frac{1}{\dim\lambda}\,\frac{D_{\lambda}^{D}(t)}{D_{\varnothing}^{D}(t)}.
\end{equation}
Moreover, for every $p,q\in\N_0$,
\begin{equation}
\label{eq:SOeven_kernel_bessel}
M_t^{D}(p,q)=I_{|p-q|}(t)+I_{p+q}(t).
\end{equation}
\end{proposition}

\begin{proof}
Insert \Cref{eq:SOeven_character_formula} and \Cref{eq:SO_even_symbol} into the general coefficient formula
\Cref{eq:one_plaq_coeff_general} and use the Weyl integration formula \Cref{eq:SOeven_weyl_integration}.
The denominator determinant in \Cref{eq:SOeven_character_formula} cancels one factor of the squared Weyl denominator,
leaving the torus integral
\begin{equation}
\label{eq:SOeven_andreief_integral}
\int_{[0,\pi]^n} W_t(\theta)\,
\det\big(\cos(A_i(\lambda)\theta_j)\big)_{1\le i,j\le n}\,
\det\big(\cos(b_i\theta_j)\big)_{1\le i,j\le n}\,d\theta_1\cdots d\theta_n.
\end{equation}
Apply Andr\'eief's identity to \Cref{eq:SOeven_andreief_integral} (with the common product measure $d\theta_1\cdots d\theta_n$)
to obtain the determinant \Cref{eq:SOeven_Dlambda_def}.
The same reduction with $\lambda=\varnothing$ gives the normalizing factor $D^D_{\varnothing}(t)$, yielding \Cref{eq:SOeven_coeff_det_ratio}.

For \Cref{eq:SOeven_kernel_bessel}, use
$\cos(p\theta)\cos(q\theta)=\tfrac12(\cos((p-q)\theta)+\cos((p+q)\theta))$ and
$\int_0^{\pi} e^{t\cos\theta}\cos(m\theta)\,d\theta=\pi I_m(t)$.
\end{proof}

\subsection{One-box inequality and telescoping}
\label{subsec:A1_Spin_even_one_box}

\begin{lemma}[One-box multiplicative bound for type $D_n$]
\label{lem:SOeven_one_box}
Fix $t>0$.
Let $\mu$ be a partition of length at most $n$ and let $i$ be an index such that $\mu+e_i$ is still a partition.
Write $A_i:=A_i(\mu)=\mu_i+b_i$.
Then
\begin{equation}
\label{eq:SOeven_one_box_bound}
\frac{a_{\mu+e_i}}{a_{\mu}}\ \le\ \frac{I_{A_i+1}(t)}{I_{A_i}(t)}.
\end{equation}
\end{lemma}

\begin{proof}
\textbf{Step 1: TN row shift bound for determinants.}
By Lemma~\ref{lem:moment_matrix_TN}(iii) the matrix $(M_t^{D}(p,q))_{p,q\ge 0}$ is TN.
Apply Lemma~\ref{lem:TN_row_shift} with the increasing column tuple $C=(0,1,\dots,n-1)$.
As before,
\begin{equation}
\label{eq:SOeven_det_shift}
\frac{D^{D}_{\mu+e_i}(t)}{D^{D}_{\mu}(t)}
\ \le\ \frac{\Delta(A(\mu+e_i))}{\Delta(A(\mu))}\,\frac{M_t^{D}(A_i+1,0)}{M_t^{D}(A_i,0)}.
\end{equation}

\textbf{Step 2: cancel the Vandermonde factor using the type-$D$ Weyl dimension ratio.}
For type $D_n$ (tensor sector) one may write
\begin{equation}
\label{eq:SOeven_dim_formula}
\dim\lambda
=\prod_{1\le i<j\le n}\frac{A_i(\lambda)^2-A_j(\lambda)^2}{b_i^2-b_j^2}.
\end{equation}
Hence, for $\mu\mapsto\mu+e_i$,
\begin{equation}
\label{eq:SOeven_dim_ratio}
\frac{\dim(\mu+e_i)}{\dim(\mu)}
=\prod_{j\ne i}\frac{(A_i+1)^2-A_j^2}{A_i^2-A_j^2}
=\prod_{j\ne i}\frac{A_i+1-A_j}{A_i-A_j}\cdot\frac{A_i+1+A_j}{A_i+A_j}.
\end{equation}
Therefore
\begin{equation}
\label{eq:SOeven_dim_cancel}
\frac{\dim(\mu)}{\dim(\mu+e_i)}\,\frac{\Delta(A(\mu+e_i))}{\Delta(A(\mu))}
=\prod_{j\ne i}\frac{A_i+A_j}{A_i+1+A_j}\ \le\ 1.
\end{equation}

\textbf{Step 3: compute the kernel entry ratio at column $0$.}
From \Cref{eq:SOeven_kernel_bessel},
\begin{equation}
\label{eq:SOeven_col0}
M_t^{D}(p,0)=I_p(t)+I_p(t)=2I_p(t).
\end{equation}
Thus
\begin{equation}
\label{eq:SOeven_Mratio}
\frac{M_t^{D}(A_i+1,0)}{M_t^{D}(A_i,0)}=\frac{I_{A_i+1}(t)}{I_{A_i}(t)}.
\end{equation}

\textbf{Step 4: combine.}
Using \Cref{eq:SOeven_coeff_det_ratio} we have
$\frac{a_{\mu+e_i}}{a_{\mu}}=\frac{\dim\mu}{\dim(\mu+e_i)}\cdot\frac{D^{D}_{\mu+e_i}}{D^{D}_{\mu}}$.
Combine \Cref{eq:SOeven_det_shift}, \Cref{eq:SOeven_dim_cancel}, and \Cref{eq:SOeven_Mratio} to obtain \Cref{eq:SOeven_one_box_bound}.
\end{proof}

\begin{corollary}[Product bound over boxes]
\label{cor:SOeven_telescoping}
Let $\lambda$ be a nonempty partition of length at most $n$.
Then
\begin{equation}
\label{eq:SOeven_product_over_boxes}
0\le a_{\lambda}(\beta,\alpha)
\ \le\ \prod_{(i,j)\in\lambda}\frac{I_{b_i+j}(t)}{I_{b_i+j-1}(t)},
\qquad t=\frac{\alpha\beta}{n},\quad b_i=n-i.
\end{equation}
\end{corollary}

\begin{proof}
Iterate \Cref{eq:SOeven_one_box_bound} along a one-box growth chain.
At the moment the $j$th box in row $i$ is added, $A_i=b_i+(j-1)$.
\end{proof}

\subsection{Crossover bound in Casimir form}
\label{subsec:A1_Spin_even_crossover}

\begin{lemma}[Casimir in partition coordinates for type $D_n$]
\label{lem:SOeven_casimir}
Let $\lambda$ be a partition of length at most $n$.
With the root-length convention of \Cref{def:metric_casimir},
\begin{equation}
\label{eq:SOeven_casimir_formula}
C_2(\lambda)=\sum_{i=1}^n \lambda_i\big(\lambda_i+2b_i\big),\qquad b_i=n-i.
\end{equation}
\end{lemma}

\begin{proof}
For type $D_n$, $\rho_i=b_i$ and $C_2(\lambda)=\langle\lambda,\lambda+2\rho\rangle$.
\end{proof}

\begin{proposition}[A1 crossover bound for $\Spin(2n)$]
\label{prop:SOeven_A1_final}
There exists an explicit constant $c_{D}(n)>0$ such that for every nontrivial $\lambda$ and every $\alpha\in[1/2,1]$, $\beta\ge 1$,
\begin{equation}
\label{eq:SOeven_A1_final}
 a_{\lambda}(\beta,\alpha)
\ \le\
\exp\Big(-c_{D}(n)\,\frac{C_2(\lambda)}{\beta+\sqrt{C_2(\lambda)}}\Big).
\end{equation}
\end{proposition}

\begin{proof}
The proof is identical to Appendix~\ref{subsec:A1_Sp_crossover} after replacing $b_i=n-i+1$ by $b_i=n-i$
(and noting $t=\alpha\beta/n$ here as well).
For completeness we record the final chain.

By \Cref{eq:SOeven_product_over_boxes} and Lemma~\ref{lem:amos_upper_bound},
\(\log a_{\lambda}\le -\sum_{(i,j)\in\lambda} \operatorname{arsinh}((b_i+j-\tfrac12)/t)\).
As in \Cref{eq:Sp_asinh_lower_sum},
\begin{equation}
\label{eq:SOeven_asinh_lower}
\sum_{(i,j)\in\lambda} \operatorname{arsinh}\Big(\frac{b_i+j-\tfrac12}{t}\Big)
\ge \frac{\sum_{(i,j)\in\lambda}(b_i+j-\tfrac12)}{t+K_{\max}},
\end{equation}
with $K_{\max}=b_1+\lambda_1-\tfrac12\le (n-1)+\lambda_1$.
Moreover
\begin{align*}
\sum_{(i,j)\in\lambda}(b_i+j-\tfrac12)
&=\frac12\sum_{i=1}^n\big(\lambda_i^2+2b_i\lambda_i\big)+\frac12\sum_{i=1}^n\lambda_i
\ge \frac12\,C_2(\lambda),
\end{align*}
using \Cref{eq:SOeven_casimir_formula} and $\lambda_i\ge 0$.
Also $\lambda_1\le \sqrt{C_2(\lambda)}$.
Combining these gives
\begin{equation}
\label{eq:SOeven_log_final}
\log a_{\lambda}\ \le\ -\frac12\,\frac{C_2(\lambda)}{t+(n-1)+\sqrt{C_2(\lambda)}}.
\end{equation}
Finally $t\le \beta/n\le \beta$ and $\beta\ge 1$ imply
$t+(n-1)+\sqrt{C_2}\le \beta+n+\sqrt{C_2}\le (n+1)(\beta+\sqrt{C_2})$.
Thus \Cref{eq:SOeven_A1_final} holds with $c_{D}(n)=1/(2(n+1))$.
\end{proof}

